\section{Introduction}
\label{sec:intro}

When a learned text representation appears to beat a structured baseline, two rival explanations are rarely separated: the representation carries genuine incremental information, or the baseline was weak and the model has latched onto stable \emph{entity identity}---in our instance, which firm filed the document---that a stronger, identity-aware reference already captures. A benchmark that credits a model against a single, uncalibrated reference cannot tell these apart, and that gap is where inflated ``representation helps'' claims live. We make this failure mode measurable in a high-stakes instance and ship a benchmark and protocol that detect and remove it: forecasting the realised volatility of S\&P~500 firms from the text of their SEC disclosures, where---under survivorship-free, leakage-free evaluation with firm-identity and maximal-baseline controls---the apparent value of disclosure text is largely a baseline-and-identity artefact.

Forecasting realised volatility---the ex-post variability of a stock's returns over a future window---is one of the most consequential prediction problems in finance, underpinning risk management, derivatives pricing, and capital regulation \citep{poon2003forecasting,andersen2003modeling}. Unlike the direction of returns, volatility is genuinely forecastable: calm and turbulent regimes cluster in time, a stylised fact exploited by the ARCH/GARCH families \citep{engle1982arch,bollerslev1986garch} and by the heterogeneous autoregressive model of realised volatility \citep[HAR-RV;][]{corsi2009har}. This strong, well-understood predictability is already extractable from the price history alone, and it sets a demanding bar that any new information source must clear.

The candidate source we examine is the text of regulatory disclosures filed with the U.S.\ Securities and Exchange Commission \citep{sec2024edgar}. A substantial literature reports that such text carries risk-relevant signal: word features of 10-K filings predict return volatility \citep{kogan2009predicting}, finance-specific sentiment dictionaries explain market reactions \citep{loughran2011liability}, and media tone moves prices \citep{tetlock2007giving}. The arrival of pre-trained Transformer encoders \citep{devlin2019bert} and long-document variants \citep{beltagy2020longformer} has renewed the expectation that richer representations might recover signal that bag-of-words methods missed. Yet this expectation sits on a fault line. Under the semi-strong efficient-markets hypothesis, public information should be impounded into prices almost immediately, leaving little residual content in week-old filings at one-to-four-week horizons. The field is also exposed to publication bias: ``text predicts markets'' is an attractive headline, and many positive results predate the strong-baseline, leakage-free protocols now considered essential for credible volatility evaluation \citep{hansen2005garch,poon2003forecasting}. Whether disclosure text adds \emph{incremental} value---predictive information beyond a strong price baseline, rather than skill merely inherited from the persistence of volatility---remains, we argue, an open question, and one that must be separated into two parts: whether the marginal information is statistically \emph{detectable}, and whether it is large enough to be \emph{realisable} as out-of-sample accuracy. Conflating the two is precisely what lets weak-baseline studies overclaim.

We answer it under deliberately rigorous conditions. We build a survivorship-free, point-in-time benchmark over S\&P~500 constituents (2010--2025), reconstructed from dated CRSP index-membership records so that firms are included exactly while they were members, including those later acquired or delisted---precisely the high-volatility exits a current-membership list would discard \citep{crsp2026data}. We then evaluate a spectrum of models---from classical bag-of-words and dictionary methods through neural encoders to price--text fusion---against price-only baselines led by HAR-RV, using strictly chronological splits (train 2010--2019, validate on the 2020--2021 COVID period, test 2022--2025) and the Diebold--Mariano test \citep{diebold1995comparing} so that every claim of superiority, parity, or inferiority is a statistical statement. Two disclosure channels are studied separately because they differ sharply in length and cadence: long-form periodic filings (10-K, 10-Q; median $\sim$29{,}278 tokens, all exceeding a 4{,}096-token window) and event-driven 8-K reports (short, frequent, timely).

This paper makes three contributions:
\begin{enumerate}
  \item \textbf{SP500Vol-Bench.} A survivorship-free, point-in-time S\&P~500 disclosure-to-volatility benchmark spanning 2010--2025: 144{,}129 SEC filings (from a 144{,}443-filing parsed corpus) aligned to 431{,}245 filing-by-horizon rows, including a successor-CIK recovery step that restores $\sim$2{,}000 filings from acquired or delisted firms. This reusable artefact is the contribution that outlasts any single result.
  \item \textbf{An incremental-value protocol, and a rigorous near-null it exposes.} We credit disclosure text only for the out-of-sample QLIKE it removes beyond a \emph{recalibrated} price forecast, adjudicated with day-clustered Diebold--Mariano inference (filings cluster in calendar time and share market shocks), Holm correction, moving-block bootstrap intervals, and a label-shuffle placebo. Under this protocol no text or fusion model beats HAR-RV out of sample: 0 of 180 standalone challenger$\times$disclosure$\times$horizon comparisons favour text. A small, placebo-confirmed increment does appear against a single recalibrated-HAR reference, on a separate 69-cell disclosure$\times$model$\times$horizon combination grid (38 of 69 cells on the seed-ensemble primary, 29 under a single seed; 0.2--5.9\% QLIKE), but it is \emph{not robust}. The protocol brackets the credit a text forecast can earn with a \emph{reference interval}: the single recalibrated HAR is the most generous reference (upper bound), a \emph{firm-identity-augmented} reference---adding the firm's own mean validation-period volatility---is the most conservative (lower bound), and a five-model maximal price pool (HAR, realised-semivariance HAR, GARCH, EGARCH, ARIMA) closes the price-side gap. Under the firm-identity reference most long-form cells flip negative, and a \emph{zero-text} firm-identity forecast reproduces the apparent ``text'' gain in 53 of 69 cells; the maximal pool absorbs the remainder; and the two surviving cell-sets are disjoint, so no cell clears both. Much of what the prior literature reads as disclosure signal is therefore baseline weakness and firm-identity confounding---controls that studies from FinText through FETILDA and \citet{kogan2009predicting} do not impose. What residual signal survives is confined to event-driven 8-K filings ($+$0.2--0.5\% QLIKE), is cross-sectional (which-firm) rather than regime-timing (when), and carries no portfolio-level economic value.
  \item \textbf{A model-spectrum comparison with a controlled prompting-vs-embedding dissociation.} Evaluated identically from bag-of-words and finance dictionaries through fine-tuned encoders (BERT, FinBERT, RoBERTa, Longformer; with a native-long-context vs.\ 512-token-truncation study) and frozen 7--8B LLM embeddings to a prompted 32-billion-parameter instruction LLM, incremental value orders as fine-tuning $\le$ TF-IDF $\le$ prompting $\le$ price: neither model scale nor domain adaptation lifts text above simple representations, let alone above the price baseline. Holding the input fixed to identical curated filing excerpts, \emph{prompting} the generative LLM extracts more of the thin residual than \emph{pooling} its same-lineage embeddings, isolating the effect to mechanism rather than input curation. Contamination controls (date-only and date-plus-ticker prompts, and an era split at the LLM's training cutoff) bound what of the residual could be memorisation, and the residual does not replicate in two further LLM families---evidence we read as bounding, not banking, the prompted-LLM signal.
\end{enumerate}

\noindent Taken together, the benchmark and its control-rich protocol form a reproducible reference point for disclosure-driven volatility forecasting---and a cautionary one: under survivorship-free, leakage-free evaluation with firm-identity and maximal-price controls, the apparent volatility signal in large-cap S\&P~500 disclosure text is largely an artefact of weak baselines and firm identity, not bankable incremental information.
